% =====================================================================
% Drift-scan feasibility for SKA-Mid-like single-dish intensity mapping
%
% SKELETON. Numbers, equations and figures are in place; the prose is not.
% Figures live in ../figures/ (gitignored). Regenerate with
%     /home/geoff/gibbs_venv_312/bin/python ../ska_hi_plots.py
% then copy the PNGs into an Overleaf figures/ folder.
% =====================================================================
% mnras.cls ALREADY loads geometry, hyperref and graphicx -- loading them again
% is the "Option clash" error. It also takes no font-size option, so no 11pt.
\documentclass[a4paper,fleqn,usenatbib]{mnras}

\usepackage{amsmath,amssymb}
\usepackage{booktabs}
\usepackage{siunitx}

% siunitx has no parsec, so \mega\parsec parses as a prefix with no unit --
% that is the "Found prefix part with no unit" error.
\DeclareSIUnit{\parsec}{pc}

\graphicspath{{figures/}{../figures/}}

% To go back to a plain article: swap the class line for
%     \documentclass[11pt,a4paper]{article}
% add   \usepackage[margin=2.4cm]{geometry} \usepackage{graphicx}
%       \usepackage[hidelinks]{hyperref}
% and change every figure* back to figure (mnras is two-column, article is not).

\newcommand{\kpar}{k_{\parallel}}
\newcommand{\kperp}{k_{\perp}}
\newcommand{\Tb}{T_{\mathrm{b}}}

% mnras wants a short running title and a short author list in the brackets.
\title[Drift scanning is information-limited]{Drift scanning for SKA-Mid-like
single-dish intensity mapping: an information-limited regime}

\author[G. Murphy et al.]{
G. Murphy$^{1}$\thanks{E-mail: TODO}
and \textit{(co-authors)}
\\
$^{1}$TODO affiliation
}

\begin{document}
\maketitle

\begin{abstract}
% TODO. Skeleton of the argument:
% (i) a parked dish cannot cross-link at any azimuth -- geometric, not a
%     matter of degree;
% (ii) the residual is set by the number of measured sky modes, not by
%      sensitivity: 17 for the drift against 34 for a constant-elevation
%      raster, and integration does not change it;
% (iii) with an injected HI signal, nothing is recoverable at the bottom of
%       Band 1 -- the post-clean residual is 347x (drift) and 24x (raster)
%       the signal, and it is the beam + prior floor, not noise;
% (iv) cross-linking improves every part of this by roughly an order of
%      magnitude, which is the practical recommendation.
\end{abstract}

% =====================================================================
\section{Introduction}
% TODO
% Context: single-dish IM, MeerKLASS, SKA-Mid Band 1. Why drift scanning is
% attractive (no moving parts, stable ground/spillover signature) and what
% the open question is.

% =====================================================================
\section{Simulations}
\label{sec:sims}

\subsection{Instrument and sky}
Single \SI{15}{\metre} dish at the Karoo site
($\phi = \ang{-30.713}$), \SI{2}{\second} sampling, \SI{3}{\hour} total
integration, HEALPix maps at $N_{\mathrm{side}} = 64$. The sky is GSM2008
rotated to equatorial coordinates. Two beams: a symmetric Gaussian with
$\mathrm{FWHM} = 1.22\lambda/D = \ang{3.99}$ at \SI{350}{\mega\hertz}, and a
tapered circular aperture ($-\SI{14}{\deci\bel}$ edge taper,
$\mathrm{FWHM} = \ang{3.84}$, first sidelobe at $\sim-\SI{23}{\deci\bel}$).

\subsection{Time-ordered data}
With no receiver temperature and no additive $T_{\mathrm{sys}}$ term, so that
the only signal is sky:
\begin{equation}
  d(t) = \bigl[ A\,s \bigr](t)\,\bigl(1 + g(t)\bigr)\bigl(1 + \eta(t)\bigr),
  \label{eq:tod}
\end{equation}
where $A$ is the pointing-and-beam operator, $g$ is $1/f$ gain drift from a
pure $f^{-2}$ flicker model, and $\eta$ is white noise of fractional variance
$\num{2.5e-6}$. Both noise terms are \emph{fractional}, so the contamination
they inject scales with sky brightness.
% TODO: note the seeding scheme -- identical realisations across scenarios is
% what makes every comparison in this paper controlled.

\subsection{Scan strategies}
Both produce exactly 5400 samples, so comparisons are at fixed clock time.
\textbf{Drift}: parked at azimuth \ang{0}, elevations \ang{52}/\ang{50}/\ang{48}
on three consecutive sidereal days. \textbf{Raster}: MeerKLASS-style
constant-elevation scan at \ang{38.19}, triangle-wave azimuth $\pm\ang{7.5}$
about \ang{45} (rising) and \ang{315} (setting), \SI{6}{\arcminute\per\second}.

\paragraph{A parked dish cannot cross-link.}
A dish at fixed $(\mathrm{az}, \mathrm{el})$ is stationary in the rotating
Earth frame, so its boresight traces a constant-declination circle swept in
$+\mathrm{RA}$. The track position angle is \ang{90} at \emph{every} azimuth
(verified numerically at az $0/45/90/270/315$); azimuth selects which
declination, not the scan direction. The two raster passes cross at \ang{73.9},
which is the one thing the drift cannot do at any pointing.

% =====================================================================
\section{Map-making}
\label{sec:mapmaking}

The map-maker is a Wiener filter. With $N$ the noise covariance, $S$ the prior
covariance and $\mu$ the prior mean,
\begin{equation}
  m = \bigl(A^{\top} N^{-1} A + S^{-1}\bigr)^{-1}
      \bigl(A^{\top} N^{-1} d + S^{-1}\mu\bigr).
  \label{eq:wiener}
\end{equation}
Equation~\eqref{eq:wiener} is \emph{affine} in $d$: a fixed linear map
$W = (A^{\top}N^{-1}A + S^{-1})^{-1}A^{\top}N^{-1}$ plus a constant prior term.
Two consequences are used throughout. An injected sky component adds linearly,
\begin{equation}
  m(d + A s) - m(d) = W A s,
  \label{eq:linearity}
\end{equation}
and solving a component's own TOD with $\mu = 0$ returns $WAs$ directly, at one
solve rather than two. Equation~\eqref{eq:linearity} is verified numerically to
$5\times10^{-6}$.

\paragraph{The beam + prior floor.}
Solving with \emph{noiseless} data leaves a residual
\begin{equation}
  r_{\mathrm{floor}} = (I - WA)\,s,
  \label{eq:floor}
\end{equation}
i.e.\ sky the strategy never measured, filled in by the prior. This term does
not integrate down, and it is the quantity that limits everything below.

\paragraph{Prior choice.}
% TODO: expand. A beam-smoothed-truth prior flatters the drift, which leans on
% it 1.40x against the raster's 1.05x. All numbers quoted here use a FLAT prior,
% constant at the patch mean, so that every spatial mode in the answer comes
% from the data.

% =====================================================================
\section{How much sky does a strategy actually measure?}
\label{sec:modes}

Eigen-analysis of $A^{\top}N^{-1}A$ counts the independently constrained sky
modes. The result is grid-independent (identical at $N_{\mathrm{side}} = 32$,
64, 128, 256):

\begin{table}[h]
\centering
\begin{tabular}{lccc}
\toprule
 & effective modes & modes $> 1\%$ of $\lambda_{\max}$ & pixels solved \\
\midrule
Drift  & 3.9 & \textbf{17} & 317 \\
Raster & 5.5 & \textbf{34} & 684 \\
\bottomrule
\end{tabular}
\caption{The strategy, not the pixelisation, sets how much sky is measurable.
Cross-linking doubles the number of measured modes.}
\label{tab:modes}
\end{table}

% TODO: the consequence -- Tsys/sqrt(N_samples) overstates the information by
% ~19x, because it assumes every pixel is independently measured.

% =====================================================================
\section{Foreground reconstruction}
\label{sec:foreground}

% TODO: the 001--005 results. Suggested content:
% - 1/f is not the limiter: the residual decomposes in quadrature to a
%   21.4 K beam+prior floor, 7.3 K white, 4.8 K 1/f. Removing 1/f entirely
%   buys 1.6%.
% - Cross-linking on a flat prior: floor -32.6% per pixel, -46.4% at beam
%   scale.
% - Integration: 3 h -> infinity moves the total only 21.9 -> 15.7 K, because
%   the asymptote IS the floor. Information-limited, not noise-limited.

% =====================================================================
\section{HI signal recovery}
\label{sec:hi}

\subsection{Injected signal and geometry}
A log-normal HI cube (Gaussian density $\rightarrow$ HI bias $\rightarrow$
log-normal $\rightarrow$ linear RSD $\rightarrow$ $\Tb(z)$) is generated at
$z_{\mathrm{eff}} = 2.80$, giving $\Tb = \SI{0.518}{\milli\kelvin}$ and
$b_{\mathrm{HI}} = 1.559$. It is placed on the HEALPix grid by a
three-dimensional lightcone interpolation: each (pixel, channel) sample is
evaluated at its true comoving position $r(z)\,\hat{n}$.

At $z \simeq 2.8$ a \SI{15}{\metre} dish resolves \SI{401}{\mega\parsec}
transverse, so $\kperp$ reaches only \SI{0.016}{\per\mega\parsec} while $\kpar$
runs \numrange{0.012}{0.192}. The accessible three-dimensional $k$-space is a
sliver near the $\kpar$ axis, and all spectra below are quoted against $\kpar$
for that reason.

\subsection{Estimator}
With $w_i$ a Blackman taper over $N_c$ channels of comoving depth $\Delta r$,
$L = N_c \Delta r$, and $\overline{w^2}$ the mean squared taper,
\begin{equation}
  P_{ab}(k) = \frac{\Delta r^{2}}{L\,\overline{w^{2}}}
    \Bigl\langle \operatorname{Re}\bigl[\tilde a(k)\,\tilde b^{*}(k)\bigr]
    \Bigr\rangle_{\mathrm{pix}},
  \qquad k = 2\pi f_{\mathrm{rfft}} .
  \label{eq:pkpar}
\end{equation}
The taper is not cosmetic: the foregrounds are $\sim10^{4}$ times the HI, so
band-edge leakage would otherwise dominate every high-$\kpar$ bin.

\subsection{Foreground cleaning and the transfer function}
Cleaning removes the leading $n$ eigenmodes of the channel--channel covariance
$C = \operatorname{cov}(x - \bar{x})$. With $U$ their eigenvectors,
\begin{equation}
  x_{\mathrm{clean}} = x - \bigl[ U U^{\top}(x - \bar{x}) + \bar{x} \bigr].
  \label{eq:pca}
\end{equation}
This is the one step that is \emph{not} linear in the data, which is why the
signal loss must be measured by injection rather than computed. Following
\citet{Cunnington2023}, with $\mathcal{C}_n$ the clean at $n$ modes and $X_m$ an
independent mock,
\begin{equation}
  T(k) = \frac{P\bigl(\mathcal{C}_n[X + X_m] - \mathcal{C}_n[X],\ X_m\bigr)}
              {P(X_m,\,X_m)} .
  \label{eq:tf}
\end{equation}
The cross-power in the numerator avoids the positive noise bias an auto-power of
the difference would carry. Because $X_m$ is each strategy's own map-made
response $WAs_m$, $T$ is a ratio within one arm and is therefore
depth-independent.

\subsection{Common resolution}
The beam narrows \SI{12.5}{\percent} across \SIrange{350}{400}{\mega\hertz}.
Each channel is reconvolved to the widest beam in the band with a Gaussian
kernel
\begin{equation}
  \theta_{k}(\nu) = \sqrt{\theta_{\max}^{2} - \theta^{2}(\nu)},
  \label{eq:commonres}
\end{equation}
exact for the Gaussian beam used here. This halves the raster's residual and
leaves the drift unchanged.

% ---------------------------------------------------------------------
\subsection{Results}

\begin{figure*}[t]
\centering
\includegraphics[width=\textwidth]{hi_transfer_function.png}
\caption{Fraction of map-made HI surviving the clean, $T(\kpar)$, per number of
modes removed. Band is $\pm1\sigma$ over 20 independent mocks. Loss concentrates
at low $\kpar$, where the modes are smoothest along the line of sight and least
distinguishable from foreground.}
\label{fig:tf}
\end{figure*}

\begin{figure}[t]
\centering
\includegraphics[width=0.72\textwidth]{hi_residual_vs_modes.png}
\caption{Post-clean residual in units of the HI it contains. Dashed is the
pipeline without the common-resolution step of Eq.~\eqref{eq:commonres}, solid
with it; both on the same interior pixels. Removing more modes lowers the
residual and costs signal, but never brings it near unity.}
\label{fig:residual}
\end{figure}

\begin{figure*}[t]
\centering
\includegraphics[width=\textwidth]{hi_ablation.png}
\caption{Noiseless (``floor'') and full data agree to \SI{0.2}{\percent} in
power -- the open rings sit on the line -- while noise alone runs
$\sim38\times$ lower. The blocker is structural, not a matter of sensitivity.}
\label{fig:ablation}
\end{figure*}

\begin{figure*}[t]
\centering
\includegraphics[width=\textwidth]{hi_frequency_structure.png}
\caption{Left: radial HI power surviving the map-maker. Right: the same
quantity split into its frequency-varying and frequency-coherent parts. The
drift retains \SI{38}{\percent} of the former while amplifying the latter
$3.14\times$; a plain rms ratio ($0.924$) hides this entirely.}
\label{fig:freqstruct}
\end{figure*}

\begin{table}[t]
\centering
\begin{tabular}{lccc}
\toprule
residual / HI, 4 modes & as run, 277\,px & interior, 119\,px & common res., 119\,px \\
\midrule
Drift  & 570.6 & 322.1 & 347.3 \\
Raster & 65.9  & 48.5  & \textbf{24.0} \\
\midrule
Cross-linking gain & $8.7\times$ & $6.6\times$ & $\mathbf{14.4\times}$ \\
\bottomrule
\end{tabular}
\caption{Two corrections, kept separable: restricting to pixels more than
\ang{3} from the patch edge removes an edge inflation, and the reconvolution of
Eq.~\eqref{eq:commonres} then halves the raster's residual while leaving the
drift unchanged.}
\label{tab:residual}
\end{table}

% =====================================================================
\section{Why a different cleaning basis cannot help}
\label{sec:rank}

Every subspace method -- PCA, ICA, GMCA, NMF, kernel PCA -- removes a rank-$n$
subspace of the channel--channel covariance and they differ only in how that
subspace is chosen. The relevant measurement is therefore how many modes each
component occupies.

\begin{table}[h]
\centering
\begin{tabular}{lcccc}
\toprule
modes for & \SI{90}{\percent} & \SI{99}{\percent} & \SI{99.9}{\percent}
  & $\lambda_{10}/\lambda_{1}$ \\
\midrule
GDSM foreground   & \textbf{1} & \textbf{1} & \textbf{1} & \num{2.6e-17} \\
Floor, drift      & 2 & 6 & 10 & \num{6.3e-4} \\
Floor, raster     & 1 & 2 & 5  & \num{5.7e-5} \\
HI, drift         & 14 & \textbf{23} & 28 & --- \\
HI, raster        & 17 & \textbf{25} & 30 & --- \\
\bottomrule
\end{tabular}
\caption{The foreground is rank one over a \SI{50}{\mega\hertz} block; the HI is
the only component that is not compressible.}
\label{tab:rank}
\end{table}

\begin{figure*}[t]
\centering
\includegraphics[width=\textwidth]{hi_rank.png}
\caption{Eigenspectrum and cumulative variance of the channel--channel
covariance for each component.}
\label{fig:rank}
\end{figure*}

\begin{figure*}[t]
\centering
\includegraphics[width=\textwidth]{hi_eigenvectors.png}
\caption{The leading eigenvectors themselves. Mode 1 of the foreground is a
single smooth power law; its modes 2--4 carry
$\lambda/\lambda_{1} \sim 10^{-17}$ and are floating-point noise.}
\label{fig:eigenvectors}
\end{figure*}

% TODO: the argument. One PCA mode removes the foreground entirely (after which
% the drift residual is still 896x the HI), so from mode 2 onward the filter is
% fighting the instrumental floor rather than foregrounds. The floor is itself
% low-rank, so a subspace method CAN capture it -- but its modes are the smooth
% ones, and the HI's low-kpar power sits in the same few. Mean cosine of the
% principal angles between the rank-6 floor and HI subspaces is 0.60 (drift)
% and 0.74 (raster). No choice of basis escapes an overlap.

% =====================================================================
\section{Discussion}
% TODO. Points worth making:
% - The scope of the negative result: drift scanning at the BOTTOM of Band 1
%   on 3 h. Foreground-to-HI contrast is 3978 at 350 MHz against 862 at 1 GHz,
%   and the beam is 2.6x more chromatic; this is close to a worst case.
% - Why more integration and more dishes do not help: the noise is already
%   ~38x below the floor, so driving it to zero improves the residual by 2%.
% - What does help: mode count, i.e. geometry. And modelling the floor rather
%   than filtering it -- Eq.~\eqref{eq:floor} is known exactly in simulation.
% - Relation to MeerKLASS-type analyses: their strongest results are
%   cross-correlations, where a residual uncorrelated with the tracer inflates
%   the variance rather than biasing the answer; and survey volume averages a
%   residual down in a way a 277-pixel patch cannot.

% =====================================================================
\section{Conclusions}
% TODO

% =====================================================================
% \bibliographystyle{mnras}   % switch on once you move to a .bib
\begin{thebibliography}{9}
\bibitem[Cunnington et al.(2023)]{Cunnington2023}
Cunnington S., et al., 2023, MNRAS % TODO complete
\bibitem[Wang et al.(2021)]{Wang2021}
Wang J., et al., 2021, MNRAS % TODO complete
\end{thebibliography}

\end{document}
